
%(BEGIN_QUESTION)
% Copyright 2011, Tony R. Kuphaldt, released under the Creative Commons Attribution License (v 1.0)
% This means you may do almost anything with this work of mine, so long as you give me proper credit

En vanlig, men uheldig trend i mange CV-er er lister over egenskaper kandidaten beskriver om seg selv. Eksempler på dette er:

\begin{itemize}
\item{} «Hardtarbeidende og motivert»
\item{} «Lett å samarbeide med»
\item{} «Kreativ problemløser»
\item{} «Dedikert»
\item{} «Velstelt og profesjonell»
\item{} «Kritisk tenker»
\item{} «Samvittighetsfull»
\item{} «Alltid punktlig»
\end{itemize}

Forklar hvorfor slike egenpåstander har liten eller ingen verdi i en CV, sett fra arbeidsgiverens perspektiv.

\vskip 50pt

Gi deretter et konkret eksempel på hvordan man kan presentere en av de positive egenskapene i en CV på en faktabasert måte (ikke subjektivt), og som gjør en bedre jobb med å fange arbeidsgiverens oppmerksomhet.

\vskip 50pt

\underbar{file i00723}
%(END_QUESTION)





%(BEGIN_ANSWER)

Et generelt svar på den andre delen av spørsmålet er å vise til konkrete prestasjoner i jobb eller skole, der den positive egenskapen var avgjørende for at man lyktes.

\vskip 10pt

{\it Merk at setningen over ikke utgjør et fullstendig svar. Du må fortsatt vise et konkret eksempel på hvordan dette kan gjøres i en CV.}

%(END_ANSWER)





%(BEGIN_NOTES)

Konkrete eksempler kan formuleres omtrent slik:

\vskip 10pt {\narrower \noindent \baselineskip5pt

«Ansvarlig for ledelse av 25 ansatte fordelt på tre skift.»

\vskip 10pt

«Kåret til “Most Valuable Employee” sommeren 2005.»

\vskip 10pt

«Ansvarlig for salg på \$135,000 i første kvartal 2003.»

\vskip 10pt

«Ledet prosjektteam i vårkvartalet ved installasjon av styringssystem i BTCs operative fiskeanlegg.»

\vskip 10pt

«Var ASB-visepresident ved BTC i studieåret 2008-2009.»

\vskip 10pt

«Organiserte tre bedriftsbesøk etter skoletid for instrumentstudentklubben i høstkvartalet 2006.»

\par} \vskip 10pt

%INDEX% Karriere, CV-skriving

%(END_NOTES)


